% Source: Work-by-gas.pdf, pages 1-2. No lecture number; placed after the Dieterici EOS, which it cites.
\subsubsection*{Example: work done by a gas in isothermal expansion}

An example of the application of more accurate equations of state (EOS): the work done by a gas as it expands
under isothermal conditions. Recall the calculation for an ideal gas with $pV = N\kB T$: the work is the area under
the curve,
\begin{equation}
  W = \int_{V_1}^{V_2}p\,\dd V = N\kB T\int_{V_1}^{V_2}\frac{\dd V}{V} = N\kB T\ln\left(\frac{V_2}{V_1}\right).
\end{equation}

\begin{diagram}
\begin{tikzpicture}[scale=0.8]
  \draw[->] (0,0) -- (4.5,0) node[below] {$V$}; \draw[->] (0,0) -- (0,3) node[left] {$p$};
  \fill[pattern=north east lines,pattern color=blue] (1.2,0) -- plot[domain=1.2:2.6] (\x,{2.2/\x}) -- (2.6,0) -- cycle;
  \draw[thick,blue,domain=0.75:4.2,samples=50] plot (\x,{2.2/\x});
  \draw[dashed,blue] (1.2,0) -- (1.2,1.83); \draw[dashed,blue] (2.6,0) -- (2.6,0.85);
  \node[below] at (1.2,0) {$V_1$}; \node[below] at (2.6,0) {$V_2$};
\end{tikzpicture}
\diagramcaption{The work done in an isothermal expansion is the area under the isotherm.}
\end{diagram}

Now do the calculation for a gas with the van der Waals EOS, with $v = V/N$:
\begin{equation}
  p = \frac{\kB T}{v-b} - \frac{a}{v^2} = \frac{\kB T}{\left(V/N\right)-b} - \frac{a}{\left(V/N\right)^2} = \frac{N\kB T}{V-bN} - \frac{aN^2}{V^2},
\end{equation}
so
\begin{align}
  W &= \int_{V_1}^{V_2}p\,\dd V = N\kB T\int_{V_1}^{V_2}\frac{\dd V}{V-bN} - aN^2\int_{V_1}^{V_2}\frac{\dd V}{V^2}\nonumber\\
    &= \left.N\kB T\ln(V-bN)\right|_{V_1}^{V_2} + \underbrace{\left.aN^2\left(\frac{1}{V}\right)\right|_{V_1}^{V_2}}_{aN^2\left(1/V_2-1/V_1\right)}\nonumber\\
    &= N\kB T\ln\left(\frac{V_2-bN}{V_1-bN}\right) + aN^2\left(\frac{V_1-V_2}{V_1V_2}\right)\nonumber\\
  &= N\kB T\ln\left[\frac{V_2\left(1-bN/V_2\right)}{V_1\left(1-bN/V_1\right)}\right] - \frac{aN^2\Delta V}{V_1V_2},
\end{align}
where $\Delta V = V_2-V_1>0$. The logarithm is
\begin{equation}
  N\kB T\left\{\ln\left(\frac{V_2}{V_1}\right) + \ln\left[\frac{1-bN/V_2}{1-bN/V_1}\right]\right\};
\end{equation}
expand for small $bN/V$:
\begin{equation}
  \ln\left[\frac{1-bN/V_2}{1-bN/V_1}\right]\simeq\ln\left[\left(1-b\frac{N}{V_2}\right)\left(1+b\frac{N}{V_1}\right)\right]
  \simeq bN\left(\frac{1}{V_1}-\frac{1}{V_2}\right) = \frac{bN(V_2-V_1)}{V_1V_2}.
\end{equation}
So
\begin{equation}
  W\simeq N\kB T\ln\left(\frac{V_2}{V_1}\right) + \underbrace{\frac{bN^2\kB T\Delta V}{V_1V_2} - \frac{aN^2\Delta V}{V_1V_2}}_{bN^2\left(\kB T-a/b\right)\Delta V/(V_1V_2) = b\kB(T-T_{\text{Boyle}})N^2\Delta V/(V_1V_2)},
\end{equation}
where here, for the vdW EOS, $\kB T_{\text{Boyle}} = a/b$. Compare with the work done by an ideal gas,
$W_{\text{ideal}} = N\kB T\ln\left(V_2/V_1\right)$:
\begin{equation}
  \Delta W = W - W_{\text{ideal}} = b\kB(T-T_{\text{Boyle}})\frac{N^2\Delta V}{V_1V_2}.
\end{equation}
N.B.: $bN = V_{\text{total excluded vol}}\equiv V_{\text{excl.}}$. So
\begin{align}
  \frac{\Delta W}{W_{\text{ideal}}} &= \frac{(bN)N\kB T\left(1-T_{\text{Boyle}}/T\right)(V_2-V_1)/(V_1V_2)}{N\kB T\ln(V_2/V_1)}\nonumber\\
  &= \frac{\left(V_{\text{excl.}}/V_1\right)\left(1-V_1/V_2\right)\left(1-T_{\text{Boyle}}/T\right)}{\ln\left(V_2/V_1\right)},
\end{align}
so
\begin{equation}
  \begin{cases}
    W>W_{\text{ideal}} & \text{if } T>T_{\text{Boyle}}\\
    W<W_{\text{ideal}} & \text{if } T<T_{\text{Boyle}} .
  \end{cases}
\end{equation}
This is clear from the definition $W = \int_{V_1}^{V_2}p\,\dd V$ and the property of the pressure:
\begin{equation}
  p = \frac{\kB T}{v-b} - \frac{a}{v^2} = \frac{\kB T\rho_n}{1-b\rho_n} - a\rho_n^2,
\end{equation}
where $\rho_n = N/V = 1/v$; for $b/v = b\rho_n\ll 1$,
\begin{equation}
  p\simeq\kB T\rho_n(1+b\rho_n) - a\rho_n^2 = \kB T\rho_n\left[1+(b-\beta a)\rho_n\right];
\end{equation}
now $p_{\text{ideal}} = N\kB T/V = \kB T\rho_n$, so
\begin{equation}
  \frac{p}{p_{\text{ideal}}} = 1 + b\left(1-\beta\frac{a}{b}\right)\rho_n = 1 + b\left(1-\frac{T_{\text{Boyle}}}{T}\right)\rho_n,
\end{equation}
so $p>p_{\text{ideal gas}}$ if $T>T_{\text{Boyle}} = a/b$ and $p<p_{\text{ideal gas}}$ if $T<T_{\text{Boyle}}$, and
the same applies for $\int_{V_1}^{V_2}p\,\dd V = W$.

Some values of $T_{\text{Boyle}}$, with $T_{\text{boil}} = T_{\text{liq.}}$ (at $p = 1$~atm) for comparison:
\begin{center}
\begin{tabular}{lcccc}
  gas & $T_{\text{boil}}$ & $T_c$ & $T_{\text{Boyle}}$ & $T_c/T_{\text{Boyle}}$\\ \hline
  Ar & 87.3 K & 151 K & $\sim 408$ K & 0.37\\
  N$_2$ & 77.35 K & 126 K & $\sim 326$ K & 0.39\\
  O$_2$ & 90.2 K & 155 K & $\sim 410$ K & 0.38\\
  CO$_2$ & (sublimates at $p = 1$ atm) & 304 K & $\sim 718$ K & 0.42
\end{tabular}
\end{center}
Reference for these values of $T_{\text{Boyle}}$: R.~Estrada-Torres et al., Fluid Phase Equilibria \textbf{258}, 148
(2007).

Recall that for the van der Waals EOS, $\kB T_c = (8/27)a/b$ and $\kB T_{\text{Boyle}} = a/b$, so
\begin{equation}
  \text{vdW}\Rightarrow\frac{T_c}{T_{\text{Boyle}}} = \frac{8}{27} = 0.296,
\end{equation}
which is slightly less than the experimental values. For the Dieterici EOS, $\kB T_c = a/(4b)$,
$\kB T_{\text{Boyle}} = a/b$, so $T_c/T_{\text{Boyle}} = 1/4 = 0.25$.
